\section{Unsymmetric spectral elements}
\label{sec:formulation}

\subsection{Problem and Petrov--Galerkin discretisation}

We consider linear plane elasticity on $\Omega\subset\R^2$,
\begin{equation}
  -\nabla\cdot\bm\sigma(\vect u) = \vect b \ \text{in } \Omega,\qquad
  \bm\sigma = \matr C\,\bm\varepsilon(\vect u),\qquad
  \vect u = \bar{\vect u}\ \text{on } \Gamma_u,\qquad
  \bm\sigma\vect n = \bar{\vect t}\ \text{on } \Gamma_t ,
  \label{eq:strong}
\end{equation}
with the strain vector $\bm\varepsilon=[\varepsilon_{xx},\varepsilon_{yy},
\gamma_{xy}]^\T$ and the plane stress or plane strain matrix $\matr C$. The mesh
consists of quadrilaterals $\Omega_e=\vect F_e(\hat\Omega)$,
$\hat\Omega=[-1,1]^2$, with an isoparametric geometry map $\vect F_e$ of degree $p$
in each reference coordinate $\bxi=(\xi,\eta)$. The element nodes are the images
$\bx_k=\vect F_e(\bxi_k)$, $k=1,\dots,n=(p+1)^2$, of the tensor-product GLL nodes
$\bxi_k$; for simplicity of notation $p_\xi=p_\eta=p$ unless stated otherwise.

The \emph{parametric} functions $N_k=\hat N_k\circ\vect F_e^{-1}$, where $\hat
N_k\in\Qp_p(\xi,\eta)$ are the Lagrange polynomials through the GLL nodes, span the
conforming test space $\mathcal W_h\subset H^1(\Omega)$. The \emph{metric} trial
functions $M_k^e$ span an element space $\mathcal U_e$ of dimension $n$ and are
interpolatory, $M_k^e(\bx_l)=\delta_{kl}$; the global trial function is
$\vect u_h|_{\Omega_e}=\sum_k \hat{\vect u}_k M_k^e$ with the nodal displacements
$\hat{\vect u}_k$ shared by adjacent elements. The trial functions are in general
discontinuous across element edges (they coincide only at the nodes). The
discrete problem reads: find the nodal values such that
\begin{equation}
  \sum_e\int_{\Omega_e}\bm\varepsilon(\vect w_h)^\T\matr C\,\bm\varepsilon(\vect u_h)\dd\Omega
  = \int_\Omega\vect w_h^\T\vect b\dd\Omega + \int_{\Gamma_t}\vect w_h^\T\bar{\vect t}\dd\Gamma
  \quad\forall\,\vect w_h\in\mathcal W_h,\ \vect w_h|_{\Gamma_u}=\vect 0 .
  \label{eq:pg}
\end{equation}
In matrix form the element stiffness is
\begin{equation}
  \matr K_e=\int_{\hat\Omega}\tilde{\matr B}_L^\T\,\matr C\,\matr B_R\dd\xi\dd\eta,
  \qquad \tilde{\matr B}_L=\det\matr J\,\matr B_L ,
  \label{eq:Ke}
\end{equation}
where $\matr B_L$ contains the Cartesian derivatives of the parametric functions and
$\matr B_R$ those of the metric functions, evaluated at $\vect F_e(\bxi)$.
$\tilde{\matr B}_L$ is assembled from $\adj\matr J\,\nabla_{\bxi}\hat N_k$, so
$\det\matr J$ does not appear in a denominator. $\matr K_e$ is unsymmetric unless
$\mathcal U_e=\mathcal W_h|_{\Omega_e}$. The consistent mass matrix is
$\matr M_e=\int_{\hat\Omega}\rho\,\matr N^\T\matr M\det\matr J\dd\xi\dd\eta$ with
the parametric ($\matr N$) and metric ($\matr M$) function values; with GLL nodal
quadrature it is diagonal, because both families are interpolatory at the GLL nodes.
We refer to the formulation as \PP\ (parametric test and trial functions, i.e.\ the
classical spectral element) or \PM\ (parametric test, metric trial functions), and to
\PM\ with the hybrid trial space of \cref{sec:hybrid} as \PH.

\subsection{Distortion immunity}

\begin{proposition}[Distortion immunity]
\label{prop:immunity}
Let every element map $\vect F_e$ be bijective, let $\Pp_q(x,y)\subset\mathcal U_e$ for
every element, let the exact solution of \eqref{eq:strong} satisfy
$\vect u\in[\Pp_q(x,y)]^2$, let the quadrature of \eqref{eq:Ke} and of the load vector be
exact for the resulting integrands, and let the global system matrix (after imposing the
Dirichlet conditions) be non-singular. Then the solution of \eqref{eq:pg} is
$\vect u_h=\vect u$, independently of the shape of the elements.
\end{proposition}
\begin{proof}
Because $\mathcal U_e$ is interpolatory and contains $\vect u|_{\Omega_e}$, the
interpolant of the exact nodal values is $\vect u$ itself on every element, so
$\bm\varepsilon(\vect u_h)=\bm\varepsilon(\vect u)$. As $\vect w_h$ is continuous,
integration by parts over $\Omega$ gives
$\int_\Omega\bm\varepsilon(\vect w_h)^\T\bm\sigma(\vect u)\dd\Omega
=\int_\Omega\vect w_h^\T\vect b\dd\Omega+\int_{\Gamma_t}\vect w_h^\T\bar{\vect t}\dd\Gamma$.
Hence the exact nodal values satisfy \eqref{eq:pg}, and they are its unique solution.
\end{proof}

\begin{remark}[Folded elements]
\label{rem:folded}
If $\vect F_e$ is not injective (``folded'' elements, $\det\matr J$ changing sign), the
parametric functions are not well defined in physical space and the argument above does
not apply. The element matrices \eqref{eq:Ke}, however, are defined in reference
coordinates with the signed $\det\matr J$ (through $\adj\matr J$), as are the loads. For
such elements exactness of polynomial solutions has been observed for low-order elements
\citep{Rajendran2010technique} and is confirmed by our patch tests
(\cref{sec:num-patch}); we do not give a proof here.
\end{remark}
This is the higher-order form of the conditions of \citet{Rajendran2010technique}.
The parametric space violates the hypothesis on curved elements: for a geometry map
of degree $g$ in each variable, $x^ay^b\circ\vect F_e\in\Qp_{g(a+b)}$, so that
$\Pp_k(x,y)\circ\vect F_e\subset\Qp_p(\xi,\eta)$ is guaranteed only for $gk\le p$
\citep[cf.][]{Arnold2002Approximation}. For bilinear maps ($g=1$), the classical
tensor-product element is complete of degree $p$; the loss of completeness under
purely angular distortion reported by \citet{Lee1993effects} concerns serendipity
spaces.

\subsection{Metric shape functions in skew coordinates}
\label{sec:skew}

Following \citet{Yuan1993New} and \citet{Xie2016Modified}, the metric functions are
written in the corotational skew coordinates
\begin{equation}
  \vect s=(s,t)^\T=\matr J_0^{-\T}(\bx-\bx_0),\qquad
  \bx_0=\vect F_e(\vect 0),\quad
  \matr J_0=\frac{\partial(x,y)}{\partial(\xi,\eta)}\Big|_{\bxi=\vect 0},
  \label{eq:skew}
\end{equation}
which require $\det\matr J_0\ne0$, with the Jacobian convention
$\matr J=\bigl[\begin{smallmatrix}x_{,\xi}&y_{,\xi}\\x_{,\eta}&y_{,\eta}\end{smallmatrix}\bigr]$.

\begin{lemma}[Affine invariance]
\label{lem:affine}
Let $\bx\mapsto\matr A\bx+\vect c$ with $\det\matr A\ne0$ be applied to all nodes of
an element. Then the skew coordinates of every material point are unchanged. In
particular, any space defined in $(s,t)$ (or in $(\xi,\eta)$) is invariant under
rigid-body motions, uniform scaling and every other affine transformation of the
element.
\end{lemma}
\begin{proof}
The transformed map is $\vect F_e'=\matr A\vect F_e+\vect c$, hence
$\matr J_0'^\T=\matr A\matr J_0^\T$ and $\bx'-\bx_0'=\matr A(\bx-\bx_0)$, so
$\vect s'=\matr J_0^{-\T}\matr A^{-1}\matr A(\bx-\bx_0)=\vect s$.
\end{proof}
Because $\vect s$ is affine in $\bx$, $\Qp_p(s,t)\supset\Pp_{p}(s,t)=\Pp_p(x,y)$, so
the space $\mathcal U_e=\Qp_p(s,t)$ satisfies the hypothesis of \cref{prop:immunity}
with $q=p$ on every element. On an affine element (parallelogram), $\vect F_e$ is
affine and $\vect s=\bxi$, so the metric and parametric functions coincide. Tensor
products of monomials in centroid-shifted Cartesian coordinates
\citep{Rajendran2003novel} span $\Qp_p(x,y)$, which contains $\Pp_p(x,y)$ as well but
depends on the orientation of the Cartesian frame (e.g.\ $xy$ versus $st$); this
makes the element frame dependent \citep{Ooi2008Remedies}.

Given a basis $\{\phi_c\}_{c=1}^n$ of $\mathcal U_e$, the metric functions are
computed from the generalised Vandermonde matrix and the basis values at the query
points $\bx_q$ (e.g.\ the images of the quadrature points),
\begin{equation}
  V_{ck}=\phi_c(\bx_k),\qquad P_{cq}=\phi_c(\bx_q),\qquad
  \matr N=\matr V^{-1}\matr P,\qquad N_{lq}=M_l^e(\bx_q),
  \label{eq:vandermonde}
\end{equation}
and analogously for the derivatives. Throughout, $\Lam_N$ denotes the Lebesgue
constant of tensor-product GLL interpolation on $\hat\Omega$, and for an
interpolatory space $\mathcal U_e$ with shape functions $M_l^e$
\begin{equation}
  \lambda_{\mathcal U}(\bx)=\sum_{l=1}^n|M_l^e(\bx)|,\qquad
  \Lam_{\mathcal U}=\max_{\bx\in\Omega_e}\lambda_{\mathcal U}(\bx)
  \label{eq:lebesgue}
\end{equation}
are its Lebesgue function and Lebesgue constant. $\Lam_{\mathcal U}=\|\mathcal
I_{\mathcal U}\|_{L^\infty\to L^\infty}$ is the norm of the nodal interpolation
operator $\mathcal I_{\mathcal U}$, so that
$\|f-\mathcal I_{\mathcal U}f\|_\infty\le(1+\Lam_{\mathcal U})\inf_{v\in\mathcal
U_e}\|f-v\|_\infty$. Note that the parametric Lebesgue constant is independent of the
element map, because $\sum_k|N_k(\vect F_e(\bxi))|=\sum_k|\hat N_k(\bxi)|$.
